\subsection{A limitation of the single-place arithmetic-source method}
\label{subsec:source-dimension-obstruction}

\begin{proposition}[A dimension obstruction to a single local arithmetic source]
\label{prop:pro3-source-dimension-obstruction}
Let $K/\mathbb Q$ be a finite Galois extension containing $\mu_3$,
let $d=[K:\mathbb Q]$, and let $t$ be the number of its primes above
three.  Let $S$ be a finite set of finite primes of $K$, containing
all primes above three and stable under $\operatorname{Gal}(K/\mathbb Q)$,
and write $s=|S|-t$.  For any $v\mid3$ one has
\[
 \dim_{\mathbb F_3}\ker\left(
 H^1(G_{K,S},\mathbb F_3)\longrightarrow H^1(K_v,\mathbb F_3)
 \right)
 \ge \frac{d(t-2)}{2t}+t+s-2.
\]
In particular, the homomorphism from $G_{K_v}$ to the maximal
pro-three quotient of $G_{K,S}$ cannot be surjective if $t\ge3$,
or if $t=2$ and $s>0$.

For every such $K$ containing $\mathbb Q(\zeta_{11})$, the second
conclusion applies when $S$ contains all primes above both three and
eleven.  For every such $K$ containing $\mathbb Q(\zeta_{13})$,
the first conclusion applies even when $S$ consists only of the
primes above three.
\end{proposition}
\begin{proof}
Since $\mu_3\subset K$, the constant coefficient module $\mathbb F_3$
can be identified with $\mu_3$.  Kummer theory embeds
$\mathcal O_{K,S}^{\times}/(\mathcal O_{K,S}^{\times})^3$
in $H^1(G_{K,S},\mu_3)$.  The field $K$ is totally imaginary,
with $r_2=d/2$ complex places.  Dirichlet's unit theorem gives
$S$-unit rank $r_2-1+|S|$; the three-primary roots of unity
contribute one further dimension modulo cubes.  Hence
\[
 \dim H^1(G_{K,S},\mathbb F_3)\ge d/2+t+s.
\]
No class-number assumption is used.

Local Kummer theory identifies the target with $K_v^\times/K_v^{\times3}$.
The valuation supplies one dimension.  A sufficiently small principal
unit group is the additive group of a finite-index $\mathbb Z_3$
lattice in $K_v$, and the three-primary roots of unity supply one
additional dimension.  Thus the target has dimension
$[K_v:\mathbb Q_3]+2=d/t+2$; the equality of local degrees uses
Galoisness.  Subtracting gives the claimed lower bound.  Surjectivity
on the pro-three arithmetic source would make restriction injective
on its homomorphisms to $\mathbb F_3$, contradicting a positive bound.

There are two primes above three in $\mathbb Q(\zeta_{11})$, since
$\operatorname{ord}_{11}(3)=5$.  Every extension containing that
field has at least two such primes, while the allowed primes above
eleven give $s>0$.  Similarly
$\operatorname{ord}_{13}(3)=3$, so $\mathbb Q(\zeta_{13})$ has
four primes above three, and every containing extension has
$t\ge4$.  This proves the final assertions.
\end{proof}

\begin{remark}[What the obstruction does not say]
\label{rem:pro3-source-dimension-scope}
This is an obstruction to the particular universal arithmetic-source
argument, not to injectivity for a specified geometric coefficient
module, and not a construction of a noncuspidal Selmer section.
A global $G_{\mathbb Q}$-equivariant cocycle that vanishes at a chosen
place also has conjugate local restrictions constrained at the other
places above three.  Joint localization or an isotypic cohomology
calculation can therefore be stronger than the single-place argument
tested here.  Compact geometric Kummer classes may also be unramified
at the exponent prime after potential good reduction, allowing a
smaller set $S$.  When $t=2$ and $s=0$, the displayed dimension bound
is zero and supplies no obstruction.
\end{remark}
